Na zijn ontwaken keek Boeddha met het oog van een Boeddha naar de wereld. Hij zag dat alle wezens in de wereld in vuur en vlam stonden door de koorts van zintuiglijk verlangen, de koorts van woede en de koorts van waanideeën. In een tekst die kort daarvoor werd gereciteerd, sprak Boeddha over hoe de ogen in brand staan vanwege passie, afkeer en illusie. Dit vuur verspreidt zich door de oren, de neus, de tong, het lichaam en de geest. Elk zintuig brandt met deze krachten. Zo zag Boeddha de wereld.

Maar als wij naar de wereld kijken, ziet alles er heel anders uit. Wij zien bepaalde beelden, geluiden, geuren, smaken en lichamelijke gewaarwordingen als iets dat heerlijk is, iets waar we van houden. En toch zegt Boeddha dat dit liefde voor vuur is. Eén van de moeilijkheden van het spirituele pad is te accepteren dat Boeddha misschien gelijk heeft, dat wij dingen verkeerd waarnemen.

Boeddha gebruikt een beeld uit de teksten: een uit leprose genezen man. Terwijl hij ziek was, had hij zijn pijnlijke wonden gestild door ze te dicht bij vuur te houden, waardoor de jeuk tijdelijk verdween. Na genezing wilde hij echter het vuur niet meer aanraken, omdat zijn waarneming nu correct was: vuur is brandend, heet en gevaarlijk. De boodschap is dat zolang we niet ontwaakt zijn, onze waarnemingen verstoord zijn. Daarbij helpt het niet dat het lichaam voedsel nodig heeft om te overleven en dat sommige smaken goed zijn terwijl andere niet lekker zijn. Het is gemakkelijk om gehecht te raken aan deze smaken en van daaruit aan andere zintuiglijke ervaringen: visuele, auditieve, zelfs mentale 'smaken'.

Het is belangrijk te reflecteren op het gevaar van zintuiglijkheid. Hierbij gaat het niet alleen om lust, maar ook om onze hartstocht voor zintuiglijke verlangens: de plannen die we maken om zintuiglijk genot te verkrijgen. Vaak genieten we meer van het bedenken van deze plannen dan van het daadwerkelijke genot. Denk bijvoorbeeld aan hoe fantaseren over voedsel vaak prettiger is dan het eten zelf. Onze echte gehechtheid zit in het denken, waar de geest het meest lijkt te branden.

Bij het najagen van zintuiglijke genoegens doen we vaak dingen die zelfzuchtig of onaangenaam zijn—allemaal om te krijgen wat we willen. En als we het eenmaal hebben, hoe lang blijft het bestaan? De Thaise monnik Ajaan Suwat benadrukte dat zintuiglijke genoegens vluchtig zijn. Vraag jezelf af: waar zijn die fijne ervaringen van gisteren of vorige week nu? Ze zijn verdwenen. Je herinnert je nog de moeite die je hebt gedaan om ze te verwerven, maar ze verdwijnen telkens zonder iets blijvends achter te laten behalve herinneringen, die zowel aangenaam als onaangenaam kunnen zijn. Soms denk je alleen aan wat er verloren is, of aan dingen die je moest doen, misschien onethisch, om ze te verkrijgen, en ervaart zelfs consequenties.

Ajaan Suwat wees ook op een intrigerend maar onthullend feit: de meeste moorden in de wereld vinden plaats tussen mensen die seks met elkaar hebben gehad. Als seks zo'n goed, opbeurend iets was, waarom leidt het dan tot zulke extreme reacties? Surrogaatliefde zoals jaloezie en wraak ontstaan vaak nadat de daad is verricht, in plaats van tevredenheid. Dit toont aan dat een deel van ons al weet dat onze slavernij aan zintuiglijke verlangens ons degradeert en niet eervol is. Toch proberen we deze bewustwording te blokkeren.

Juist dit deel van de geest moeten we onderzoeken als we enig inzicht, concentratie of vrijheid willen bereiken. Daarom richt zoveel van de Boeddhistische Canon zich op de nadelen van zintuiglijkheid. Boeddha vergelijkt deze gehechtheid met het dragen van een brandende fakkel tegen de wind: als je hem niet loslaat, verbrand je. Hij beschrijft het ook als een druppel honing op een mes, of een roofvogel met een stukje vlees, terwijl andere roofvogels proberen het weg te rukken. Er is een lange lijst met beelden die de gevaren illustreren.

Mensen vragen vaak waarom Boeddha zo'n nadruk legt op de negatieve kanten van zintuiglijkheid. Het antwoord is dat we zó verslaafd zijn aan het zogenaamd positieve, dat we constante herinneringen nodig hebben aan de andere kant. Hierdoor kan het deel van ons dat inziet dat er nadelen zijn, weten dat het niet alleen is. In een wereld waar bijna iedereen geobsedeerd is door zintuiglijkheid, voelt dit deel van de geest zich vaak geïsoleerd.

Wat Boeddha probeert te doen is dat isolement doorbreken. Hij voedt dat deel van de geest dat een vorm van geluk verlangt die niet afhankelijk is van zintuiglijke genoegens. Door hierover te reflecteren, zal de geest meer geneigd zijn om goedheid en welzijn te vinden op een andere manier—een vorm die de teksten omschrijven als 'losgekoppeld van zintuiglijkheid'.

Dit proces helpt niet alleen om negatieve zintuiglijke verlangens los te laten; het is ook een benadering van concentratie-ontwikkeling, gericht op hoger plezier. De geest moet een alternatieve bron van voldoening vinden om zichzelf los te maken van het lagere genot. Dit hoger genot kan voortkomen uit het volledig bewonen van je eigen lichaam, terwijl je de ademhaling en de energie in je lichaam kalmeert en richt.

Boeddha benadrukt dat zonder het ervaren van plezier dat voortkomt uit afzondering van verlangens, onthechting onmogelijk is. Het noodzakelijke hogere genot wordt gecreëerd door het lichaam te laten versmelten met de ademhaling en innerlijke rust te vinden. Dit maakt het gemakkelijker om zintuiglijke verlangens los te laten.

Boeddha roept op tot een reflectie die verder gaat dan oppervlakkig verlangen. Terwijl moderne zintuiglijkheid energie verspilt en ons oplegt, biedt Boeddha een alternatief: diep innerlijke concentratie en welzijn. Dit is geen gewelddadige verwijdering van zintuigen, maar juist een fijne harmonie van innerlijke lichamelijke sensaties en mentale abstractie. Zo kan de geest aan vrijheid voorbij het zintuiglijke ontsnappen.

Waarom zijn we verslaafd aan sensualiteit? (Het antwoord van de Boeddha)

Er bestaan vijf strengen van sensualiteit. Welke vijf? Vormen die via het oog waarneembaar zijn—aangenaam, plezierig, charmant, innemend, verlangens opwekkend en verleidelijk; geluiden waarneembaar via het oor; aroma’s waarneembaar via de neus; smaken waarneembaar via de tong; en tastbare sensaties waarneembaar via het lichaam—alle aangenaam, plezierig, charmant, innemend, verlangens opwekkend en verleidelijk. Maar deze dingen zijn niet de sensualiteit zelf. Ze worden in de discipline van de nobele mensen strengen van sensualiteit genoemd. Passie voor zijn voornemens is de sensualiteit van een mens, niet de mooie sensuele genoegens die in de wereld worden gevonden. Passie voor zijn voornemens is de sensualiteit van een mens. De schoonheid blijft zoals ze is in de wereld, terwijl wijzen hun verlangen in die zin bedwingen.

En wat is de oorzaak waarom sensualiteit in werking treedt? Contact. Wat is de diversiteit in sensualiteit? Sensualiteit met betrekking tot vormen is één ding, met betrekking tot geluiden is het een andere, met betrekking tot aroma’s weer een andere, en met betrekking tot smaken wederom een andere, evenals sensualiteit met betrekking tot tastbare sensaties. En wat is het gevolg van sensualiteit? Iemand die naar sensualiteit verlangt creëert een corresponderende zelf-toestand (attabhāva) aan de kant van verdienste of tekortkoming. En wat is de beëindiging van sensualiteit? Door de beëindiging van contact treedt de beëindiging van sensualiteit in; en precies dit nobele achtvoudige pad—juiste zienswijze, juiste bedoeling, juiste spraak, juiste handelen, juiste levensonderhoud, juiste inspanning, juiste mindfulness, juiste concentratie—is de weg die leidt tot de beëindiging van sensualiteit.

Zoals gesteld in de Nibbedhika-soetra (A N 6:63):

**De onwetendheid en de pijnlijke cyclus van sensualiteit**

Wanneer een persoon geteisterd wordt door een pijnlijke gewaarwording, geniet de onwetende, doorsneemens van sensualiteit. Waarom is dat? Omdat hij geen ontsnapping uit pijnlijke gewaarwordingen kent behalve uit sensualiteit. Dit wordt duidelijk uitgelegd in de Sallattha Sutta (SN 36:6).

**De metafoor van de melaatse**

“Māgaṇḍiya, stel dat er een melaatse was, bedekt met zweren en infecties, verteerd door wormen, die schilfers van de openingen van zijn wonden afkrabde met zijn nagels, terwijl hij zijn lichaam dicht bij een put van gloeiende kolen verhitte. Zijn vrienden, metgezellen en familieleden zouden hem naar een dokter brengen. De dokter zou medicijnen bereiden, en dankzij de medicijnen zou hij genezen van zijn lepra. Nu gezond en gelukkig, vrij en meester van zichzelf, kon hij gaan waar hij wilde. Stel je nu voor dat twee sterke mannen hem mee zouden slepen naar een dergelijke put van gloeiende kolen, terwijl hij nu gezond was. Wat denk je? Zou hij zich niet verzetten?”

Māgaṇḍiya antwoordt dat de melaatse zich zou verzetten omdat het vuur heet en pijnlijk is. De Boeddha vervolgt: “Was het vuur niet altijd al pijnlijk, zowel vroeger als nu, maar waren de percepties van de melaatse toen geïmpaird, wat hem deed denken dat het vuur ‘aangenaam’ was?”

In dezelfde lijn legt de Boeddha uit dat sensualiteit altijd pijnlijk is. Alleen wanneer men is verteerd door verlangen en passie, worden de zintuigen vertroebeld, waardoor men het pijnlijke als plezierig ervaart. Zoals een melaatse troost en bevrediging vindt in het krabben van zijn wonden, zo vinden mensen een vluchtige vorm van plezier in het voldoen aan hun sensualiteit.

De maakt mensen afhankelijk. Hoe meer men zich overgeeft aan sensualiteit, des te meer het verlangen groeit en de interne ‘koorts’ van het verlangen toeneemt.

De Boeddha maakt duidelijk dat zelfs machtige koningen of ministers die temidden van alle vormen van sensueel genot leven, zonder het verlangen los te laten of die koorts van verlangen te onderdrukken, nooit vrijheid van dorst of innerlijke rust zullen kennen. Alleen degenen die echt hebben gerealiseerd wat sensualiteit omvat—haar oorsprong, verdwijnen, aantrekkingskracht, gevaar en ontsnapping—kunnen deze rust daadwerkelijk bereiken.

Thema: Lust als Sleutel tot Zelfinzicht en Godsbewustzijn
Kernpunten over seksualiteit:
Metafoor van Pandora's Box: Seksualiteit lijkt klein, maar heeft grote impact op iemands leven.

Seksualiteit overstijgt de daad zelf en weerspiegelt diepere lagen van een persoon.

Seksualiteit en creativiteit:
Seksuele energie is scheppende energie.

Seksualiteit staat symbool voor wat iemand in zijn leven wil creëren.

Verbindt met diepste verlangens en angsten.

Seksualiteit toont de meest intense wensen en behoeften van een persoon.

Verbinding seksualiteit met trauma en verlangen:
Diepste verlangens komen voort uit diepste trauma’s.

Seksuele voorkeuren onthullen verborgen trauma's en onvervulde behoeften.

Seksualiteit is dus niet oppervlakkig, maar een spiegel van kernwonden.

Seksualiteit als weg naar ontwaken:
Door seksualiteit te onderzoeken, kom je oog in oog met diepste kwetsbaarheid en echte verlangens.

Seksualiteit kan leiden tot spiritueel ontwaken.

Emotionele component van seksualiteit is cruciaal voor diep inzicht.
